\documentclass[12pt]{article}
\usepackage[utf8]{inputenc}
\usepackage[T1]{fontenc}
\usepackage[margin=1in]{geometry}
\usepackage{setspace}
\usepackage{xcolor}
\usepackage[colorlinks=true, urlcolor=blue!60!black, linkcolor=blue!60!black]{hyperref}
\setlength{\parindent}{2em}
\setlength{\parskip}{1.25em}
\renewcommand{\baselinestretch}{1.0}
\begin{document}
\begin{center}
{\LARGE \textit{NLRB v. Jones \& Laughlin Steel Corp.}, 301 U.S. 1 (1937)\par}
\vspace{0.8em}
{\large Case Brief --- Four-Page Edition\par}
\end{center}
\vspace{1em}

\section{Memory Jogger}

Even though manufacturing is intrastate when viewed in isolation, Congress may regulate labor relations at a vertically integrated steel plant under the Commerce Clause whenever industrial strife there would have a ``close and substantial relation to interstate commerce'' (\href{https://supreme.justia.com/cases/federal/us/301/1/}{301 U.S. 1}).

\section{Detailed Case Facts}

The National Labor Relations Act of 1935 (the Wagner Act) protected employees' rights to organize and bargain collectively and barred ``unfair labor practices'' ``affecting commerce,'' including coercing employees and discriminating against them to discourage union membership.

Jones \& Laughlin, the fourth-largest U.S. steel producer, ran a completely integrated enterprise: ore mines in Michigan and Minnesota, coal and limestone properties, Great Lakes steamships, terminal railroads, warehouses and fabricating shops in several states, and sales offices in twenty cities. About 75\% of its steel was shipped out of Pennsylvania. At its Aliquippa plant, which employed about 10,000 workers, it discharged ten union leaders. The NLRB found the discharges were for union activity and ordered reinstatement with back pay; the company, which had withdrawn from the hearing after contesting the Board's jurisdiction, refused to comply.

\section{Procedural History}

The Fifth Circuit denied the Board's enforcement petition, holding that the order regulated local manufacturing beyond federal power. The Supreme Court granted certiorari, heard the case with four companion ``Labor Board Cases'' on February 10--11, 1937, and reversed on April 12, 1937.

\section{Judicial Votes}

Reversed, \textbf{5--4}.
\begin{itemize}
\item \textbf{Majority:} Chief Justice Hughes, joined by Justices Brandeis, Stone, Roberts, and Cardozo.
\item \textbf{Concurrences:} None.
\item \textbf{Dissent:} Justice McReynolds, joined by Justices Van Devanter, Sutherland, and Butler.
\end{itemize}

\section{Holding}

The NLRA is constitutional as applied, and the Board's order is sustained. Congress may regulate intrastate activities, including manufacturing labor relations, that have such a close and substantial relation to interstate commerce that their control is essential or appropriate to protect that commerce from burdens and obstructions. The Act does not violate Fifth Amendment due process, provides adequate judicial review, and its reinstatement-with-back-pay remedy is not a suit at common law triggering the Seventh Amendment. Hughes's warning that federal power may not reach effects ``so indirect and remote'' as to ``obliterate the distinction between what is national and what is local'' was dicta, but it later anchored \textit{United States v. Lopez} (1995).

\section{Analysis of Opinions}

\textbf{Hughes (majority).} Invoking avoidance---``The cardinal principle of statutory construction is to save, and not to destroy''---Hughes read the Act to reach only practices ``affecting'' commerce, case by case. Employees' right to organize is a ``fundamental right,'' since a single worker dependent on daily wages is helpless against arbitrary treatment. Declining to rest on the ``stream of commerce'' metaphor of \textit{Stafford v. Wallace}, he held that ``It is the effect upon commerce, not the source of the injury, which is the criterion.'' Drawing on the \textit{Shreveport Rate Case} and Sherman Act decisions, and recalling the 1919--1920 steel strike, he found that a stoppage would be ``immediate'' and might be ``catastrophic.'' He distinguished \textit{Schechter} (remote local effects) and \textit{Carter Coal} (delegation and due process flaws). The Act did not compel agreements or bar legitimate discharges.

\textbf{McReynolds (dissenting).} For the ``Four Horsemen,'' manufacturing is distinct from commerce and subject to state police power. Firing ten workers had an effect ``indirect and remote in the highest degree,'' and the majority's theory ``would extend it into almost every field of human industry,'' destroying the federal balance and liberty of contract.

\section{Examples: Future Applications}

\begin{enumerate}
\item \textbf{Same-Side (Valid): Garment Manufacturer.} A Virginia suit maker imports most of its fabric and ships over 80\% of its suits interstate. As in \textit{Friedman-Harry Marks Clothing}, industrial disruption in a firm integrated into interstate markets burdens commerce.
\item \textbf{Same-Side (Valid): Intrastate Utility.} A New York utility powers interstate railroads and telegraph facilities. A strike would paralyze instrumentalities of commerce (as in \textit{Consolidated Edison}, 1938).
\item \textbf{Opposite-Side (Invalid): Guns Near Schools.} A statute bans possessing a gun near a school, with no jurisdictional hook. Hughes's ``so indirect and remote'' limit supports invalidation, as in \textit{Lopez}.
\item \textbf{Opposite-Side (Invalid): Custody Standards.} Congress sets custody rules for state courts because family instability lowers productivity. The causal chain is too attenuated and would yield ``a completely centralized government.''
\item \textbf{Fence-Sitter: Neighborhood Bakery.} Congress regulates a single-location bakery selling only locally. Alone its effect seems remote, but the aggregation that \textit{Jones \& Laughlin} set in motion, culminating in \textit{Wickard v. Filburn}, could sustain regulation.
\end{enumerate}

\section{Critique}

\textbf{Originalist and textualist views.} Richard Epstein argues that ``commerce'' originally meant trade, not manufacturing, and that the ``substantial effects'' test erased enumerated powers; Randy Barnett adds that Hughes's promised limit gave way to \textit{Darby} and \textit{Wickard}. Robert Pushaw counters that the founders understood commerce to include market-oriented production.

\textbf{Progressive and historical views.} Progressives praise Hughes's economic realism over the direct--indirect and production--commerce dichotomies. Historians debate whether 1937 was ``the switch in time that saved nine,'' a constitutional moment in Bruce Ackerman's sense, or, as Barry Cushman argues, an organic doctrinal evolution.

\section{Key Quotations}

\begin{quote}
``Although activities may be intrastate in character when separately considered, if they have such a close and substantial relation to interstate commerce that their control is essential or appropriate to protect that commerce from burdens and obstructions, Congress cannot be denied the power to exercise that control.''
\end{quote}
\noindent --- Chief Justice Hughes, Opinion of the Court, 301 U.S. at 37.

\begin{quote}
``The authority of the federal government may not be pushed to such an extreme as to destroy the distinction, which the commerce clause itself establishes, between commerce `among the several States' and the internal concerns of a State. That distinction between what is national and what is local in the activities of commerce is vital to the maintenance of our federal system.''
\end{quote}
\noindent --- Chief Justice Hughes, Opinion of the Court, 301 U.S. at 30.

\begin{quote}
``We have often said that interstate commerce itself is a practical conception. It is equally true that interferences with that commerce must be appraised by a judgment that does not ignore actual experience.''
\end{quote}
\noindent --- Chief Justice Hughes, Opinion of the Court, 301 U.S. at 41--42.

\end{document}
